\section{Contoh Protokol: TLS 1.3 Handshake}

\subsection{Flow Lengkap TLS 1.3}

TLS 1.3 menyederhanakan handshake dibandingkan TLS 1.2, mengurangi latensi menjadi 1-RTT (Round Trip Time). Berikut adalah alur lengkap koneksi aman antara Client (Browser) dan Server (Website).

\begin{figure}[htbp]
  \centering
  \begin{tikzpicture}[node distance=3cm, auto, thick]
    % Nodes
    \node (client) {\textbf{Client}};
    \node (server) [right=6cm of client] {\textbf{Server}};
    
    % Timeline lines
    \draw (client) -- ++(0,-10);
    \draw (server) -- ++(0,-10);
    
    % Messages
    % 1. Client Hello
    \draw[->] ($(client)+(0,-1)$) -- node[above, align=left] {\textbf{Client Hello}\\ + Key Share (e.g., ECDH public key)\\ + Cipher Suites Supported\\ + Random Nonce} ($(server)+(0,-2)$);
    
    % 2. Server Hello
    \draw[<-] ($(client)+(0,-4)$) -- node[above, align=right] {\textbf{Server Hello}\\ + Key Share (e.g., ECDH public key)\\ + Cipher Suite Selected\\ + Server Random} ($(server)+(0,-3)$);
    
    % Encryption starts here!
    \draw[dashed, red] ($(client)+(0,-4.5)$) -- node[midway, fill=white] {Enkripsi Dimulai (Handshake Keys)} ($(server)+(0,-4.5)$);
    
    % 3. Encrypted Extensions
    \draw[<-] ($(client)+(0,-5.5)$) -- node[above, align=right] {\textbf{\{Encrypted Extensions\}}\\ \textbf{\{Certificate\}}\\ \textbf{\{Certificate Verify\}} (Signature)\\ \textbf{\{Finished\}} (MAC)} ($(server)+(0,-4.5)$);
    
    % 4. Finished
    \draw[->] ($(client)+(0,-6.5)$) -- node[above, align=left] {\textbf{\{Finished\}} (MAC)} ($(server)+(0,-7.5)$);
    
    % Encryption Application Data
    \draw[dashed, blue] ($(client)+(0,-8)$) -- node[midway, fill=white] {Application Data Encryption (Traffic Keys)} ($(server)+(0,-8)$);
    
    % 5. HTTP Data
    \draw[<->, double] ($(client)+(0,-9)$) -- node[above] {\textbf{\{HTTP Request / Response\}}} ($(server)+(0,-9)$);
    
  \end{tikzpicture}
  \caption{Alur TLS 1.3 Handshake yang menunjukkan pertukaran kunci Diffie-Hellman dalam paket Hello pertama.}
\end{figure}

\subsection{Penjelasan Tahapan}

\textbf{1. Client Hello}:
\begin{itemize}
  \item Client mengirimkan \textbf{Key Share} (bagian publik dari Diffie-Hellman, missal $g^a \bmod p$) secara spekulatif.
  \item Mengirim daftar cipher yang didukung (misal: TLS\_AES\_256\_GCM\_SHA384).
  \item Mengirim \textbf{Random Nonce} untuk mencegah replay attack.
\end{itemize}

\textbf{2. Server Hello}:
\begin{itemize}
  \item Server memilih cipher suite.
  \item Server mengirimkan \textbf{Key Share} miliknya ($g^b \bmod p$).
  \item \textbf{Saat ini juga}, kedua pihak sudah bisa menghitung \textbf{Handshake Traffic Key} ($g^{ab} \bmod p$). Semua pesan berikutnya dienkripsi!
\end{itemize}

\textbf{3. Otentikasi Server}:
\begin{itemize}
  \item Server mengirim \textbf{Certificate} (mengandung Public Key RSA/ECC Server).
  \item Server mengirim \textbf{Certificate Verify}: Tanda tangan digital atas semua pesan handshake sejauh ini, menggunakan Private Key Server.
  \item Client memverifikasi tanda tangan menggunakan Public Key dari sertifikat. Ini membuktikan Server memiliki Private Key yang sesuai.
\end{itemize}

\textbf{4. Finished}:
\begin{itemize}
  \item Kedua pihak mengirim pesan \textbf{Finished} yang berisi MAC (HMAC) dari seluruh transkrip handshake. Ini memastikan integritas: tidak ada penyerang yang memodifikasi pesan Hello (misalnya downgrade attack).
\end{itemize}

\subsection{Kunci yang Diturunkan}

Dari Handshake, TLS 1.3 menurunkan beberapa kunci berbeda menggunakan HKDF (HMAC-based Key Derivation Function):

\begin{enumerate}
  \item \textbf{Handshake Traffic Keys}: Mengenkripsi sisa handshake.
  \item \textbf{Application Traffic Keys}: Mengenkripsi data HTTP.
  \item \textbf{Exporter Master Secret}: Untuk penggunaan eksternal (jarang dipakai).
  \item \textbf{Resumption Master Secret}: Untuk fast resume koneksi di masa depan (0-RTT).
\end{enumerate}

\begin{konsep}
Perbedaan utama TLS 1.3 dan 1.2 adalah \textbf{Key Exchange} dilakukan seawal mungkin. Di TLS 1.2, Key Exchange terjadi setelah Certificate exchange, membutuhkan 2 kali round-trip sebelum data dikirim. TLS 1.3 melakukannya di pesan pertama, hanya butuh 1 round-trip.
\end{konsep}
